% ============ 第5章 §5.3 Schwartz 分布 (原书 pp.108-110) ============

\section{Schwartz 分布}\label{sec:5.3}
\index{S@Schwartz 分布(空间)}




Schwartz 函数由性质良好的函数构成, 所以对象狭窄. 为了扩大考虑的对象, 考虑对偶空间是有用的想法. 这里, 将对偶空间作为拓扑对偶线性空间来考虑, 考察其性质.

$\mathcal{S}'(\mathbb{R}^{n})$ 作为拓扑对偶, 即作为 $\mathcal{S}(\mathbb{R}^{n})$ 的连续线性泛函全体构成的集合来定义.

\begin{definition}[Schwartz 分布 (空间)]\label{def:5.2}
$\mathcal{S}'(\mathbb{R}^{n})$ 是从 $\mathcal{S}(\mathbb{R}^{n})$ 到 $\mathbb{C}$ 的连续复线性算子全体构成的集合. 即,
\[
\begin{aligned}
\mathcal{S}'(\mathbb{R}^{n}) \equiv \{f &: \mathcal{S}(\mathbb{R}^{n}) \to \mathbb{C} : \langle f, \varphi + \psi \rangle = \langle f, \varphi \rangle + \langle f, \psi \rangle, \\
&\quad \langle f, a \cdot \varphi \rangle = a \cdot \langle f, \varphi \rangle, f \text{ 连续}\}.
\end{aligned}
\]
这里 a 遍历 $\mathbb{C}$ 全体, $\varphi, \psi$ 遍历 $\mathcal{S}(\mathbb{R}^{n})$ 全体. f 在 $\varphi$ 处的值 $f(\varphi)$ 按惯例记为 $\langle f, \varphi \rangle$.
\end{definition}

``小''集合的对偶一般地``大'', 如下面的例子所示, $\mathcal{S}'(\mathbb{R}^{n})$ 包含通常函数以外的东西, 实际上非常大.

\begin{example}[Dirac 分布 (Dirac delta)]\label{exa:5.8}
对 $a \in \mathbb{R}^{n}$, 代入映射 $\delta_{a}: \varphi \in \mathcal{S}(\mathbb{R}^{n}) \mapsto \varphi(a) \in \mathbb{C}$ 属于 $\mathcal{S}'(\mathbb{R}^{n})$. 实际上, 由于 $|\langle\delta_{a}, \varphi\rangle| = |\varphi(a)| \leq p_{0}(\varphi) = \sup_{x \in \mathbb{R}^{n}}|\varphi(x)|$, 从 $\mathcal{S}(\mathbb{R}^{n})$ 到 $\mathbb{C}$ 的连续性得到保证. $\delta_{a}$ 在物理学等中也是重要的. 当 $a = 0$ 时, 记为 $\delta_{a} = \delta$.
\end{example}

\begin{example}\label{exa:5.9}
对 $\varphi \in \mathcal{S}(\mathbb{R}^{n})$, $f \in L^{p}(\mathbb{R}^{n})$, $1 \leq p \leq \infty$, 由 Hölder 不等式,
\[
\int_{\mathbb{R}^{n}}|f(x)\varphi(x)|\dd x \leq \|f\|_{L^{p}}\|\varphi\|_{L^{p'}} \lesssim \|f\|_{L^{p}}\sup_{x \in \mathbb{R}^{n}}(1+|x|)^{n+1}|\varphi(x)|
\]
所以映射
\[
F_{f}: f \in L^{p}(\mathbb{R}^{n}) \mapsto \int_{\mathbb{R}^{n}}f(x)\varphi(x)\dd x \in \mathbb{C}
\]
是 $\mathcal{S}'(\mathbb{R}^{n})$ 的元. 如此, 通过由 f 给出 $F_{f}$ 生成 $\mathcal{S}'(\mathbb{R}^{n})$ 的元, 但重要的是, 由 Lebesgue 微分定理 (参照注意~\ref{note:3.2}), 这个对应是单射. 实际上, 若 $F_{f} = 0$, 则
\[
\int_{\mathbb{R}^{n}}f(y)\varphi(y)\dd y = F_{f}(\varphi) = 0.
\]
由于 $\varphi \in \mathcal{S}(\mathbb{R}^{n})$ 是任意的, 取 $\varphi_{t}(x - \cdot)$ 作为 $\varphi$,
\[
f * \varphi_{t}(x) = \int_{\mathbb{R}^{n}}f(y)\varphi_{t}(x-y)\dd y = 0.
\]
因此, 由注意~\ref{note:3.2}, 对几乎所有 $x \in \mathbb{R}^{n}$, 得到 $f(x) = \lim_{t \downarrow 0}f * \varphi_{t}(x) = 0$. 由以上的经过, $F_{f}$ 可以记为 f.
\end{example}

虽然不是所有的元都具有 $F_{f}$ 的形式, 但相当多的元可以用 $F_{f}$ 的形式表示, 所以常用与函数 f 相同的字母来表示 $\mathcal{S}'(\mathbb{R}^{n})$ 的元. $\mathcal{S}'(\mathbb{R}^{n})$ 的下面的特征刻画具体且容易理解.

\begin{theorem}[$\mathcal{S}'(\mathbb{R}^{n})$ 的特征刻画]\label{thm:5.3}
设 F 是从 $\mathcal{S}(\mathbb{R}^{n})$ 到 $\mathbb{C}$ 的 $\mathbb{C}$-线性映射. 此时, $F \in \mathcal{S}'(\mathbb{R}^{n})$ 的必要充分条件是, 存在某个 $N \in \mathbb{N}$, 使得对所有 $\varphi \in \mathcal{S}(\mathbb{R}^{n})$,
\begin{equation}\tag{5.5}
|F(\varphi)| \leq Np_{N}(\varphi)
\label{eq:5.5}
\end{equation}
成立.
\end{theorem}

\begin{proof}
式~\eqref{eq:5.5} 是充分条件. 实际上, 当 $\mathcal{S}(\mathbb{R}^{n})$ 的列 $\{\varphi_{j}\}_{j=1}^{\infty}$ 收敛于 $\varphi \in \mathcal{S}(\mathbb{R}^{n})$ 时使用式~\eqref{eq:5.5},
\[
|F(\varphi_{j}) - F(\varphi)| = |F(\varphi_{j} - \varphi)| \leq Np_{N}(\varphi_{j} - \varphi)
\]
成立, 所以当 $j \to \infty$ 时得到 $F(\varphi_{j}) \to F(\varphi)$. 下面证明必要性. 假设 $F: \mathcal{S}(\mathbb{R}^{n}) \to \mathbb{C}$ 连续, 证明满足式~\eqref{eq:5.5} 的 $N \in \mathbb{N}$ 的存在性. 由 $F(0) = 0$ 与 F 的连续性, 注意
\begin{equation}\tag{5.6}
F^{-1}(\{z \in \mathbb{C} : |z| < 1\}) = \{\varphi \in \mathcal{S}(\mathbb{R}^{n}) : |F(\varphi)| < 1\}
\label{eq:5.6}
\end{equation}
是包含 0 的 $\mathcal{S}(\mathbb{R}^{n})$ 的开集. 因此, 由例~\ref{exa:5.6}, 取 $L \in \mathbb{N}$ 足够大, 使得
\[
\{\varphi \in \mathcal{S}(\mathbb{R}^{n}) : Lp_{L}(\varphi) < 1\} \subset \{\varphi \in \mathcal{S}(\mathbb{R}^{n}) : |F(\varphi)| < 1\}.
\]
因此, 若 $\varphi \in \mathcal{S}(\mathbb{R}^{n})$ 满足 $Lp_{2L}(\varphi) = \frac{1}{2}$, 则 $|F(\varphi)| \leq 1$.

那么, 现在任意给定 $\varphi \in \mathcal{S}(\mathbb{R}^{n}) \setminus \{0\}$, $\psi \equiv \frac{1}{2Lp_{2L}(\varphi)}\varphi$ 满足 $Lp_{2L}(\psi) = \frac{1}{2}$, 所以得到 $|F(\psi)| \leq 1$. 由 $\psi$ 的定义式, 对 $\varphi \in \mathcal{S}(\mathbb{R}^{n}) \setminus \{0\}$, $|F(\varphi)| \leq 2Lp_{2L}(\varphi)$ 成立. 该式对 $\varphi = 0$ 也当然成立. 因此, 取 $N = 2L$, 就取到了满足式~\eqref{eq:5.5} 的 N, 所以式~\eqref{eq:5.5} 成立.
\end{proof}

$\mathcal{S}(\mathbb{R}^{n})$ 中引入了距离拓扑, 也给 $\mathcal{S}'(\mathbb{R}^{n})$ 引入拓扑. 一般地, 给定集合 X 的子集族 Z 时, 考虑包含 Z 的开集族全体 V, 其交集
\[
\bigcap_{\mathcal{W} \in \mathcal{V}}\{U \in 2^{X} : U \in \mathcal{W}\}
\]
构成开集族, 回忆这被称为由 Z 生成的拓扑.

\begin{definition}\label{def:5.3}
设 $\mathcal{S}'(\mathbb{R}^{n})$ 中的拓扑是由下面的集合族 Z 生成的拓扑:
\[
\mathcal{Z} = \{g \in \mathcal{S}'(\mathbb{R}^{n}) : |\langle f - g, \varphi\rangle| < \varepsilon\}.
\]
这里考虑满足 $f \in \mathcal{S}'(\mathbb{R}^{n})$, $\varphi \in \mathcal{S}(\mathbb{R}^{n})$, $\varepsilon > 0$ 的所有 f, $\varphi$, $\varepsilon$.
\end{definition}

换句话说, $\mathcal{S}'(\mathbb{R}^{n})$ 的拓扑是使映射 $F \in \mathcal{S}'(\mathbb{R}^{n}) \mapsto F(\varphi) \in \mathbb{C}$ 对所有 $\varphi \in \mathcal{S}(\mathbb{R}^{n})$ 连续的最弱拓扑.

\begin{example}\label{exa:5.10}
对 $f \in \mathcal{S}(\mathbb{R}^{n})$ 的 $\partial_{x_{j}}f$ 也可以在 $\mathcal{S}(\mathbb{R}^{n})$ 的拓扑中考虑. 即, 设 $e_{j}$ 为第 j 个基本向量,
\[
\partial_{x_{j}}f = \lim_{h \to 0}\frac{1}{h}(f(\cdot + h\mathbf{e}_{j}) - f)
\]
在 $\mathcal{S}(\mathbb{R}^{n})$ 的拓扑中成立. 实际上, 需要证明的是对所有 $\alpha, \beta \in \mathbb{N}_{0}^{n}$,
\[
\lim_{h \to 0}p_{\alpha,\beta}\left(\frac{1}{h}(f(\cdot + h\mathbf{e}_{j}) - f) - \partial_{x_{j}}f\right) = 0,
\]
但由于
\begin{align*}
&\lim_{h \to 0}p_{\alpha,\beta}\left(\frac{1}{h}(f(\cdot + h\mathbf{e}_{j}) - f) - \partial_{x_{j}}f\right) \\
&= \lim_{h \to 0}p_{\alpha,0}\left(\frac{1}{h}(\partial^{\beta}f(\cdot + h\mathbf{e}_{j}) - \partial^{\beta}f) - \partial^{\beta}\partial_{x_{j}}f\right)
\end{align*}
所以可以设 $\beta = 0$. 由于 $\beta = 0$, 可以用 $p_{N}$ 代替 $p_{\alpha,0}$ 来考虑. 由于不需要考虑偏导数, 考虑
\[
q_{N}(\varphi) = \sup_{x \in \mathbb{R}^{n}}(1+|x|)^{N}|\varphi(x)|.
\]
此时, 由中值定理
\[
F(t+s) = F(t) + sF'(t) + \frac{s^{2}}{2}F''(t+\theta s), \theta \in (0, 1)
\]
可以取到依赖于 x, h, j 等的 $t \in (0, 1)$, 使得
\[
\frac{1}{h}(f(x + h\mathbf{e}_{j}) - f(x)) - \partial_{x_{j}}f(x) = \frac{h}{2}\partial_{x_{j}}^{2}f(x + ht\mathbf{e}_{j}).
\]
一般地, 对 $y_{1}, y_{2} \in \mathbb{R}^{n}$,
\[
1 + |y_{1} + y_{2}| \leq 1 + |y_{1}| + |y_{2}| \leq (1+|y_{1}|)(1+|y_{2}|)
\]
所以设 $N \in \mathbb{N}$, $|h| < 1$,
\[
(1+|x|)^{N} \leq (1+|x + ht\mathbf{e}_{j}|)^{N}(1+|ht\mathbf{e}_{j}|)^{N} \leq 2^{N}(1+|x + ht\mathbf{e}_{j}|)^{N}.
\]
这里, 右边的不等式中使用了
\[
|h| < 1, t \in (0, 1) \Longrightarrow 1 + |ht\mathbf{e}_{j}| \leq 1 + |h| < 2.
\]
由
\[
q_{N}(\varphi) = \sup_{x \in \mathbb{R}^{n}}(1+|x|)^{N}|\varphi(x)|
\]
可知, 若 $|h| < 1$, 则
\begin{align*}
&q_{N}\left(\frac{1}{h}(f(\cdot + h\mathbf{e}_{j}) - f) - \partial_{x_{j}}f\right) \\
&= \sup_{x \in \mathbb{R}^{n}}(1+|x|)^{N}\left|\frac{1}{h}(f(x + h\mathbf{e}_{j}) - f(x)) - \partial_{x_{j}}f(x)\right| \\
&= \sup_{x \in \mathbb{R}^{n}}(1+|x|)^{N}\left|\frac{h}{2}\partial_{x_{j}}^{2}f(x + ht\mathbf{e}_{j})\right| \\
&\leq 2^{N}\sup_{x \in \mathbb{R}^{n}}(1+|x + ht\mathbf{e}_{j}|)^{N}\left|\frac{h}{2}\partial_{x_{j}}^{2}f(x + ht\mathbf{e}_{j})\right| \\
&= 2^{N-1}|h|\sup_{x \in \mathbb{R}^{n}}(1+|x + ht\mathbf{e}_{j}|)^{N}|\partial_{x_{j}}^{2}f(x + ht\mathbf{e}_{j})| \\
&= 2^{N-1}|h|\sup_{z \in \mathbb{R}^{n}}(1+|z|)^{N}|\partial_{x_{j}}^{2}f(z)|
\end{align*}
即,
\[
q_{N}\left(\frac{1}{h}(\partial^{\beta}f(\cdot + h\mathbf{e}_{j}) - \partial^{\beta}f) - \partial^{\beta}\partial_{x_{j}}f\right) = O(|h|)
\]
证明结束.
\end{example}

\begin{exercise}\label{exer:5.6}
设 $\mathcal{O} = 2^{\mathcal{S}'(\mathbb{R}^{n})}$. 证明: O 是使对所有 $\varphi \in \mathcal{S}(\mathbb{R}^{n})$ 的映射 $F \in \mathcal{S}'(\mathbb{R}^{n}) \mapsto F(\varphi) \in \mathbb{C}$ 连续的 $\mathcal{S}'(\mathbb{R}^{n})$ 中的最强拓扑.
\end{exercise}
